% =====================================================================================================
% Online Appendix C: Monte Carlo design and results (module m6_montecarlo).
% Writer: conclusion writer, 2026-09-24. Tables generated by code/paper/make_appendix_cd_tables.py
% from output/tables/m6_montecarlo_*.csv. Numbers and sources: paper/notes/writer_conclusion.md
% =====================================================================================================

\section{Monte Carlo Design and Results}\label{app:mc}
\setlength{\emergencystretch}{2em}% allow looser lines in the online appendix

This appendix documents the two Monte Carlo designs behind Sections~\ref{sec:ident} and \ref{sec:obs}.
Design~A asks what alternative experimental designs identify when they use the same total compute. Design~B
simulates an industry of labs that choose compute and allocations, and asks which production-function estimators
recover the technology from the models they release. Everything is simulated. The truth is the technology of
\citet{besiroglu2024chinchilla}: $E=1.8172$, $A=482.01$, $B=2085.43$, $\alpha=0.3478$, and $\beta=0.3658$, so that
$a=0.513$, $\gamma=0.178$, $\sigma^*=0.737$, and the compute-optimal ratio at $10^{24}$ FLOP is $M^*=18.1$ tokens
per parameter ($\ln M^*=2.898$). ``Bias'' is the Monte Carlo mean error and ``RMSE'' the root-mean-squared error
around the truth.

\subsection{Design A: What Alternative Designs Identify}\label{app:mc-designA}

\emph{Setup.}---Every design has 90 runs, ten at each of Chinchilla's nine IsoFLOP budgets ($6\times10^{18}$ to
$3\times10^{21}$ FLOP), and hence the same total training compute, $5.1\times10^{22}$ FLOP. The designs differ only
in how the runs are placed along their isocost lines (Table~\ref{tab:app-mc-designs}): exactly or nearly on the
expansion path, as a compute-optimal ladder would be; spread by IsoFLOP profiles or a factorial grid, as a designed
experiment would be; spread by optimization errors of standard deviation $s$ in $\ln(D/N)$ at given compute; or
tilted systematically by labs that follow Kaplan's allocation rule. Log loss carries multiplicative noise with
standard deviation 0.0075, the residual standard deviation of the Besiroglu fit on the 240-run Chinchilla
extraction. Because the residuals are heavy-tailed (their MAD-based standard deviation is 0.0049) and bundle
seed noise, digitization error, and misspecification, this value is an upper bound on pure seed noise. Noise is
realized after inputs are chosen, the case of \citet{zellner1966specification}.

We compare three estimators. The primal estimator is the Huber ($\delta=10^{-3}$) fit of $\ln L$ used by
\citet{hoffmann2022training} and \citet{besiroglu2024chinchilla}, minimized from nine starting values inside a
parameter box; one starting value is the truth, which favors near-truth solutions on a flat ridge and makes the
on-path results conservative. The dual estimator reads the path slope $a$ from the optima (or IsoFLOP minima) and
$\gamma$ from the frontier $L^*(C)$, and recovers $\alpha=\gamma/a$ and $\beta=\gamma/(1-a)$; it identifies the
exponents only through optimality and the Chinchilla outer exponent $\kappa=1$. The third is a profile likelihood
for $\sigma^*$ in the generalized family $L=E+(AN^{-a_1}+BD^{-b_1})^{\kappa}$ with $\kappa$ free, on an 18-point
grid over $[0.50,0.95]$, with the likelihood ratio measured against the unrestricted minimum. Its starting values
include the member of the observationally equivalent family of Proposition~\ref{prop:fdep} that reproduces the
on-path losses to $10^{-15}$; without them, the optimizer produced spurious rejections at the grid edge.

\emph{On-path data do not identify the technology through the primal estimator.}---When every run is exactly
compute-optimal ($s=0$), 51 percent of fits end on a corner with one power-law term switched off, and the RMSE of
$\hat\sigma^*$ is 0.28, against 0.004 for an IsoFLOP design spanning $\pm16\times$ the optimal size and 0.005 for
the factorial grid at the same compute (Table~\ref{tab:app-mc-designs}, Panel~A). The failure is a ridge with
boundary solutions, not merely imprecision: at $s=0$ the median error of $\hat\sigma^*$ is $-0.30$, while at
$s=0.02$ the typical fit is close (median absolute error 0.037) but the 5--95 percent range is $[0.34,0.77]$. The
normalized condition number of the information matrix rises from $1.4\times10^{3}$ in the designed experiments to
$1.5\times10^{7}$ on the path and is singular at $s=0$ (Panel~B). The compute-optimal ratio is the more demanding
object once optimality is not imposed: the RMSE of $\ln\hat M^*(10^{24})$ falls from 24 at $s=0$ to 4.5 at
$s=0.1$, 2.0 at $s=0.3$, and 0.59 at $s=1$, against 0.16 under the IsoFLOP design. At $s=0.3$ the 5--95 percent
range of $\ln\hat M^*$ is $[-0.40,6.34]$, a factor of about 850 in $M^*$. These patterns are the Monte Carlo
counterpart of Proposition~\ref{prop:info}: information about the location of the expansion path grows with the
dispersion of log wedges, information about its curvature only with their fourth power.

\emph{With $\kappa=1$ imposed, apparent precision at intermediate $s$ is functional form.}---For
$0.1\le s\le0.3$ the primal RMSE of $\hat\sigma^*$ (0.013--0.025) is no better than that of the dual estimator
(0.012), which uses only the path. The reason is analytic: with $\kappa=1$, $\alpha+\beta=\gamma/[a(1-a)]$, and
$a(1-a)$ is flat near $a=1/2$ (0.2498 at the truth), so $\sigma^*=2/(2+\gamma/[a(1-a)])\approx2/(2+4\gamma)$ is
pinned down by the on-path frontier elasticity. Freeing $\kappa$ removes that information. The 95 percent
profile interval for $\sigma^*$ is the entire grid in 88 percent of on-path replications, in 89 to 97 percent of
replications for $s\le0.2$, in 65 percent at $s=0.3$, and in 77 percent under Kaplan beliefs
(Figure~\ref{fig:app-mc-profile}). It becomes informative at $s=1$ (median width 0.104, median interval
$[0.69,0.79]$) and in the designed experiments (widths 0.016 for the IsoFLOP $\pm16\times$ design, 0.019 for the
factorial grid, and 0.047 for IsoFLOP $\pm4\times$), and it covers the truth in 93 to 100 percent of replications
in every cell. What counts as ``near-optimal'' depends on $s$ relative to the noise and on the number of runs: at
$s=0.3$ the share of flat profiles is 38, 65, and 90 percent at noise standard deviations of 0.005, 0.0075, and
0.015. Released industry models are far off the path (Llama 3 8B has $M\approx1{,}900$ against $M^*\approx20$), so
this problem concerns compute-optimal ladders and sweeps, not the cross-section of deployed models.

\emph{The dual estimator is precise on the path but inherits any systematic allocation error.}---On the path its
RMSEs are 0.0005 for $a$, 0.011 for $\sigma^*$, and 0.009 for $\ln M^*$. When labs optimize against Kaplan's
beliefs ($N\propto C^{0.73}$), the same estimator is biased by $+0.217$ in $a$, $-0.088$ in $\sigma^*$, and
$-3.92$ in $\ln M^*$, understating $M^*$ about 50-fold; the primal estimator on the same data is roughly unbiased
but weakly identified (RMSE of $\hat\sigma^*$ 0.10).

\emph{On the path, standard inference is either wrong or uninformative.}---The Wald covariance is numerically
singular or the fit is at a corner in 35 percent of on-path replications. The pairs bootstrap as usually
implemented, with every draw warm-started at the full-sample estimate, covers $\sigma^*$ in only 67 percent of
replications. Restarting every draw from all nine starting values restores coverage (92 percent in a separate
check with 60 replications) but yields a median interval width of 0.43, against 0.016 under the IsoFLOP design.
With designed variation all methods agree: Wald and bootstrap coverage is 95 to 97 percent for the IsoFLOP $\pm16\times$
and factorial designs and at $s=1$. Digitizing the losses as in the Epoch extraction barely matters for an IsoFLOP
design (RMSE of $\hat\sigma^*$ 0.0041 against 0.0036).

% Online Appendix Table C1 (v3: allocation-error scale written varsigma; start-set table dropped). Writer appendices.
% Round 3 (WP3, P14(n)): factorial grid not on the nine budgets; start region stated as within Hoffmann et al.'s grid.
\begin{table}[tp]
\centering
\caption{Monte Carlo Design A: What Alternative Designs Identify at Equal Compute}
\label{tab:app-mc-designs}
\footnotesize
\setlength{\tabcolsep}{2.8pt}
\begin{tabular}{@{}lcccccc@{}}
\toprule
 & \multicolumn{4}{c}{Primal estimator ($\kappa=1$)} & \multicolumn{2}{c}{Dual estimator} \\
\cmidrule(lr){2-5}\cmidrule(l){6-7}
Design & $a$ & $\gamma$ & $\sigma^*$ & $\ln M^*(10^{24})$ & $\sigma^*$ & $\ln M^*(10^{24})$ \\
\midrule
\multicolumn{7}{@{}l}{\textit{Panel A. Bias [root-mean-squared error]}} \\[1pt]
(i) On path, $\varsigma=0.02$ & 0.050 & 0.058 & $-$0.166 & $-$1.08 & 0.000 & 0.00 \\
 & [0.296] & [0.096] & [0.222] & [23.8] & [0.011] & [0.01] \\[1pt]
(ii) IsoFLOP, $\pm16\times$ & 0.001 & 0.000 & 0.000 & $-$0.02 & 0.000 & 0.02 \\
 & [0.007] & [0.003] & [0.004] & [0.16] & [0.016] & [0.14] \\[1pt]
(ii$'$) IsoFLOP, $\pm4\times$ & 0.000 & 0.000 & 0.000 & 0.00 & $-$0.001 & 0.02 \\
 & [0.015] & [0.008] & [0.009] & [0.31] & [0.016] & [0.36] \\[1pt]
(iii) Factorial grid & $-$0.001 & 0.000 & 0.000 & 0.02 & -- & -- \\
 & [0.012] & [0.005] & [0.005] & [0.28] & -- & -- \\[1pt]
(iv) Optimizing, $\varsigma=0$ & 0.140 & 0.093 & $-$0.240 & $-$11.8 & 0.000 & 0.00 \\
 & [0.331] & [0.113] & [0.272] & [27.7] & [0.012] & [0.00] \\[1pt]
(iv) Optimizing, $\varsigma=0.1$ & 0.014 & $-$0.032 & $-$0.006 & $-$0.25 & 0.000 & 0.00 \\
 & [0.225] & [0.051] & [0.025] & [4.58] & [0.012] & [0.05] \\[1pt]
(iv) Optimizing, $\varsigma=0.3$ & 0.000 & $-$0.005 & $-$0.002 & 0.00 & 0.000 & $-$0.01 \\
 & [0.098] & [0.016] & [0.013] & [1.97] & [0.012] & [0.14] \\[1pt]
(iv) Optimizing, $\varsigma=1$ & $-$0.002 & 0.000 & $-$0.001 & 0.04 & $-$0.001 & $-$0.04 \\
 & [0.029] & [0.011] & [0.012] & [0.59] & [0.015] & [0.48] \\[1pt]
(v) Kaplan beliefs & $-$0.028 & $-$0.028 & $-$0.023 & 0.34 & $-$0.088 & $-$3.92 \\
 & [0.265] & [0.072] & [0.102] & [4.97] & [0.088] & [3.92] \\[1pt]
\midrule
\multicolumn{7}{@{}l}{\textit{Panel B. Inference for $\sigma^*$ and the profile likelihood with $\kappa$ free}} \\[1pt]
 & Corner & Wald & Bootstrap & \multicolumn{3}{c}{95\% profile set, $\kappa$ free} \\
\cmidrule(l){5-7}
 & fits (\%) & coverage (\%) & coverage (\%) & Median bounds & Flat (\%) & Flat$'$ (\%) \\
\cmidrule(l){2-7}
(i) On path, $\varsigma=0.02$ & 72 & 100$^{\dagger}$ & 44$^{\S}$ & [0.05, 0.99] & 91 & 94 \\
(ii) IsoFLOP, $\pm16\times$ & 0 & 96 & 97 & [0.729, 0.745] & 0 & 0 \\
(ii$'$) IsoFLOP, $\pm4\times$ & 0 & 94 & -- & [0.714, 0.760] & 0 & 0 \\
(iii) Factorial grid & 0 & 95 & 97 & [0.727, 0.746] & 0 & 0 \\
(iv) Optimizing, $\varsigma=0$ & 87 & 100$^{\dagger}$ & -- & [0.05, 0.99] & 95 & 97 \\
(iv) Optimizing, $\varsigma=0.1$ & 0 & 99$^{\dagger}$ & -- & [0.08, 0.99] & 21 & 93 \\
(iv) Optimizing, $\varsigma=0.3$ & 0 & 93 & 96$^{\S}$ & [0.42, 0.99] & 0 & 65 \\
(iv) Optimizing, $\varsigma=1$ & 0 & 95 & 95 & [0.69, 0.79] & 0 & 0 \\
(v) Kaplan beliefs & 0 & 92$^{\dagger}$ & -- & [0.08, 0.99] & 21 & 84 \\
\bottomrule
\end{tabular}
\par\smallskip
\begin{tablenotes}Designs, truth ($a=0.513$, $\gamma=0.178$, $\sigma^*=0.737$, $\ln M^*(10^{24})=2.90$) and noise as described in the text: 90 runs with equal total compute, on Chinchilla's nine budgets except in the factorial grid, 500 replications per design; (v) labs follow Kaplan's rule $N\propto C^{0.73}$ with errors $\varsigma=0.1$. Primal: Huber ($\delta=10^{-3}$) fit of $\ln L$ from nine starting values drawn at random within the region of \citet{hoffmann2022training}'s initialization grid, none at the truth, inside a parameter box. Dual: $a$ from the (noisy) optima or IsoFLOP minima, $\gamma$ from the frontier (eight starting values on a grid), $\alpha=\gamma/a$, $\beta=\gamma/(1-a)$; not computed for the factorial grid. Panel B: corner = estimate on the box boundary or one power-law term below 2 percent of reducible loss at every design point; Wald = LAD-sandwich delta-method 95 percent interval ($^{\dagger}$: covariance singular or corner in more than 5 percent of replications, 79 percent on the path; coverage among computable ones); bootstrap = pairs-bootstrap percentile interval (149 draws, first 150 replications), each draw warm-started at the replication's estimate ($^{\S}$: in a separate check with 60 replications and 49 draws, restarting every draw also from the nine starting values moves on-path coverage from 33 to 78 percent and the median width from 0.04 to 0.43; at $\varsigma=0.3$ both give 100 percent); 95 percent profile set = likelihood-ratio confidence set for $\sigma^*$ in $L=E+(AN^{-a_1}+BD^{-b_1})^{\kappa}$ on a 29-point grid over $[0.05,0.99]$, with the likelihood ratio measured against the unrestricted optimum (first 150 replications): medians of the lower and upper bounds, share of replications whose set is the whole grid (flat) and share whose set contains all of $[0.50,0.95]$, the range searched in the first version (flat$'$). Monte Carlo standard errors of coverage rates near 95 percent are 1.0, 1.8 and 2.8 percentage points with 500, 150 and 60 replications.\end{tablenotes}
\begin{tablenotes}[Source]Author's simulations; replication code \texttt{code/analysis/m6\_montecarlo}.\end{tablenotes}
\end{table}


\begin{figure}[tp]
\centering
\includegraphics[width=\textwidth]{m6_montecarlo_profile_ci}
\caption{Profile-Likelihood Intervals for $\sigma^*$ with the Outer Exponent Free}\label{fig:app-mc-profile}
\begin{figurenotes}
Design~A, first 150 replications per design. Left: median width of the 95 percent profile-likelihood interval for
$\sigma^*$ in $L=E+(AN^{-a_1}+BD^{-b_1})^{\kappa}$ with $\kappa$ free, on the grid $[0.50,0.95]$ (maximum width
0.45), against the standard deviation $s$ of allocation errors for optimizing labs; horizontal lines give the
IsoFLOP $\pm16\times$ and factorial designs at the same compute. Right: share of replications whose interval is
the whole grid. Log-loss noise standard deviation 0.0075; 90 runs per design.
\end{figurenotes}
\end{figure}

\subsection{Design B: A Simulated AI Industry}\label{app:mc-designB}

\emph{Setup.}---Forty labs each produce six generations of models. In each generation a lab releases one to four
models whose compute budgets are roughly $8\times$ apart. Productivity is Hicks-neutral and equals
$\omega_{ft}=0.15(t-1)+x_{ft}$, where the lab component follows $x_{ft}=0.8\,x_{f,t-1}+\xi_{ft}$ with stationary
standard deviation 0.25; seed noise has standard deviation 0.05 in $y=-\ln(L-E)$. The dispersion of $\omega$ is assumed, chosen to match
the 41-fold within-developer 90--10 ratio in effective compute of \citet{mertens2026secret}; TFP growth of 0.15
per generation is also assumed. A lab-generation's log budget is $c_{ft}=\bar c_t+h_f-p_{ft}+\nu_{ft}$, where
$\bar c_t$ grows by $\ln3$ per generation from $10^{21}$ FLOP, $h_f$ is lab size (standard deviation 0.7),
$p_{ft}$ is an observed compute-price shifter (0.4), and $\nu_{ft}$ is noise (0.4). Four compute rules map
productivity into budgets: (a) exogenous budgets; (b) funding that responds to current productivity
($+2x_{ft}$); (c) a loss target that the lab reaches by buying compute; and (d) compute provisioned one
generation ahead ($+2x_{f,t-1}$), the timing assumption of \citet{ackerberg2015identification}. Labs minimize
lifetime compute given inference demand $T=3D^*(c)\theta$ with $\ln\theta\sim N(0,0.8^2)$, so that the mean log
wedge is 0.46, and make allocation errors with standard deviation 0.2 in $\ln(D/N)$. Variants add release
selection (about 18 percent of trained models unreleased), classical measurement error in $\ln D$ (standard
deviation 0.15), data-augmenting lab heterogeneity, and a flagship-only sample. The estimators, applied to released
models, are listed in Table~\ref{tab:app-mc-industry}; $E$ is known to all of them except the pooled Huber fit that
mimics current practice in machine learning.

\emph{Transmission bias is diluted but not removed by common fixes.}---Pooled nonlinear least squares with
generation effects overstates $\gamma$ by 0.023 (13 percent) under funding and 0.018 under predetermined budgets,
and understates it by 0.051 (29 percent) under a loss target (Table~\ref{tab:app-mc-industry}, Panel~A;
Figure~\ref{fig:app-mc-industry}). With one
model per lab-generation under funding the bias is $+0.119$, matching the closed form of
Proposition~\ref{prop:transmission} applied to this design (0.118). Lab fixed effects remove about two-thirds of
the bias under funding and four-fifths under predetermined budgets, but not under the target rule ($-0.021$), and
almost none of it with flagships only ($+0.113$), because $\omega$ is persistent rather than fixed. The bias loads
on the frontier elasticity, not on the allocation: the regression of $\ln N$ on $\ln C$ recovers $a$ with
$|\text{bias}|\le0.003$ in every scenario (Panel~B, IV row).

\emph{Within-family variation is the observational analog of an IsoFLOP experiment.}---Lab $\times$ generation
fixed effects, which use only the size menu within a lab-generation, are unbiased for $\gamma$, $a$, and TFP
growth under every compute rule ($|\text{bias}|\le0.002$), because siblings share $\omega$ and their tier and
inference-demand differences are exogenous. The estimator needs at least two released siblings, is unavailable
with flagships only, and is hurt by measurement error in $D$ (bias in $\ln M^*$ of $-0.33$ at an error standard
deviation of 0.15, similar to lab fixed effects, $-0.28$). Estimators in the spirit of
\citet{ackerberg2015identification} remove the transmission bias when the timing assumption is right
($|\text{bias}\ \gamma|\le0.002$ with lagged instruments only) and fail when it is wrong ($-0.026$ under the
target rule if current compute is treated as predetermined). Both variants, however, use the current
within-family deviations of inputs as instruments, without which the lagged moment set is under-identified, so they
borrow the same within-family variation; the lagged-instrument version is also five times less precise for $a$
(RMSE 0.12 to 0.13 against 0.026). The simulation therefore does not show that proxy or timing estimators work on flagship-only panels.

\emph{First-order-condition estimators mistake inference demand for technology.}---Stacking the within-family
loss equation with the training-only first-order condition ($w=1$) leaves the exponents essentially unbiased but
biases $\ln M^*$ by $+1.25$ to $+1.28$ in every scenario (Panel~D), overstating $M^*$ about 3.5-fold. The closed
form is $2\,\mathrm{E}[\ln w]/(\alpha+\beta)=2\times0.457/0.714=1.28$: the first-order-condition residual is the
inference wedge of Proposition~\ref{prop:wedge}, and ignoring it makes over-trained models look like evidence
that data are cheap. A noisy usage proxy for $T$ (log noise 0.3) removes the bias ($-0.02$ to $-0.03$, RMSE at
most 0.05). The instrumental-variable estimator that reads $M^*$ off the observed path inherits the same
contamination ($+1.27$ to $+1.32$), although its median bias for $\gamma$ is small ($-0.002$ to $+0.005$; median
absolute error 0.027 to 0.031) when releases are not selected.

\emph{Current practice absorbs more than half of algorithmic progress into the scaling law.}---The pooled Huber
fit with $E$ free and no time effects, the specification common in the machine-learning literature, understates
TFP growth by 0.081 to 0.090 of a true 0.150 per generation under rules (a), (b), and (d) (Panel~C), and
overstates $\ln M^*$ by 0.56 to 0.68. Holding $E$ at the truth isolates the scale-attribution channel of
\citet{nagy2013statistical}: the same pooled regression still misses 36 to 43 percent of TFP growth and steepens
the frontier by 0.050 to 0.059. Roughly two-thirds of the absorbed progress is attributed to scale and one-third to
a downward-biased asymptote $\hat E$.

\emph{Selection and its correction interact with transmission.}---With exogenous budgets and a notability
threshold for release, pooled estimates understate $\gamma$ by 0.008 and overstate TFP growth by 0.011; a
\citet{heckman1979sample} correction, which assumes that the inputs of unreleased models are observed and that a
release-policy shifter is excluded, halves both. Under predetermined budgets, selection ($-$) and transmission
($+$) offset (net $+0.011$), and the correction raises the bias to $+0.018$ by removing only the selection part.
Selection also invalidates the compute-price instrument (median bias $-0.022$ and $-0.033$). The industry
calibration is illustrative, not estimated; only signs and rankings should be read from Design~B, and the
magnitudes scale with the assumed ratio of productivity-driven to total variance in compute.

\begin{table}[tp]
\centering
\caption{Monte Carlo Design B: Production-Function Estimators in a Simulated AI Industry}
\label{tab:app-mc-industry}
\footnotesize
\setlength{\tabcolsep}{2.5pt}
\begin{tabular}{@{}lcccc@{}}
\toprule
Compute rule: & (a) Exogenous & (b) Funding & (c) Target & (d) Predetermined \\
\midrule
\multicolumn{5}{@{}l}{\textit{Panel A. Frontier elasticity $\gamma$ (truth 0.178)}} \\
Pooled Huber, $E$ free (ML practice) & 0.019 [0.031] & 0.027 [0.036] & $-$0.004 [0.022] & 0.025 [0.035] \\
Pooled NLS + generation FE & $-$0.001 [0.010] & 0.023 [0.024] & $-$0.053 [0.056] & 0.018 [0.019] \\
Lab FE & 0.000 [0.006] & 0.008 [0.009] & $-$0.022 [0.024] & 0.004 [0.007] \\
Lab $\times$ generation FE & 0.000 [0.002] & 0.000 [0.003] & 0.000 [0.002] & 0.000 [0.002] \\
ACF, current $c$ as instrument & 0.000 [0.010] & 0.007 [0.009] & $-$0.027 [0.034] & 0.000 [0.010] \\
ACF, lagged instruments only & $-$0.001 [0.011] & $-$0.001 [0.011] & $-$0.002 [0.011] & $-$0.001 [0.011] \\
IV (compute price) + path$^{\ddagger}$ & $-$0.002 [0.031] & 0.005 [0.031] & $-$0.002 [0.027] & $-$0.002 [0.028] \\
FOC system, $T$ ignored & $-$0.002 [0.003] & $-$0.002 [0.003] & $-$0.003 [0.003] & $-$0.002 [0.003] \\
FOC system, $T$ proxied & $-$0.001 [0.001] & $-$0.001 [0.001] & $-$0.001 [0.001] & $-$0.001 [0.001] \\
\addlinespace
\multicolumn{5}{@{}l}{\textit{Panel B. Allocation exponent $a$ (truth 0.513)}} \\
Pooled Huber, $E$ free (ML practice) & $-$0.006 [0.111] & $-$0.004 [0.115] & $-$0.004 [0.082] & $-$0.010 [0.108] \\
Pooled NLS + generation FE & $-$0.002 [0.085] & 0.004 [0.042] & $-$0.060 [0.191] & 0.000 [0.045] \\
Lab FE & 0.001 [0.066] & 0.001 [0.043] & $-$0.010 [0.112] & 0.001 [0.060] \\
Lab $\times$ generation FE & 0.000 [0.027] & 0.001 [0.028] & 0.000 [0.025] & 0.002 [0.027] \\
ACF, current $c$ as instrument & 0.004 [0.117] & 0.001 [0.054] & $-$0.055 [0.241] & 0.009 [0.121] \\
ACF, lagged instruments only & $-$0.001 [0.126] & $-$0.001 [0.121] & 0.000 [0.126] & 0.003 [0.130] \\
IV (compute price) + path$^{\ddagger}$ & 0.000 [0.003] & 0.000 [0.002] & $-$0.003 [0.004] & 0.000 [0.002] \\
FOC system, $T$ ignored & $-$0.002 [0.004] & $-$0.002 [0.004] & $-$0.004 [0.005] & $-$0.002 [0.004] \\
FOC system, $T$ proxied & $-$0.001 [0.003] & $-$0.001 [0.003] & $-$0.001 [0.003] & $-$0.001 [0.003] \\
\addlinespace
\multicolumn{5}{@{}l}{\textit{Panel C. TFP growth per generation (truth 0.150)}} \\
Pooled Huber, $E$ free (ML practice) & $-$0.081 [0.081] & $-$0.090 [0.090] & $-$0.036 [0.039] & $-$0.088 [0.088] \\
Pooled NLS + generation FE & 0.002 [0.012] & $-$0.024 [0.027] & 0.054 [0.055] & $-$0.018 [0.021] \\
Lab FE & 0.001 [0.010] & $-$0.009 [0.013] & 0.023 [0.025] & $-$0.003 [0.010] \\
Lab $\times$ generation FE & 0.001 [0.009] & 0.000 [0.009] & 0.000 [0.009] & 0.001 [0.009] \\
ACF, current $c$ as instrument & 0.002 [0.010] & $-$0.007 [0.012] & 0.024 [0.026] & 0.002 [0.010] \\
ACF, lagged instruments only & 0.003 [0.010] & 0.002 [0.009] & 0.002 [0.009] & 0.003 [0.010] \\
IV (compute price) + path$^{\ddagger}$ & 0.005 [0.035] & $-$0.003 [0.032] & 0.001 [0.032] & 0.003 [0.032] \\
FOC system, $T$ ignored & 0.003 [0.009] & 0.002 [0.010] & 0.003 [0.009] & 0.004 [0.010] \\
FOC system, $T$ proxied & 0.001 [0.009] & 0.001 [0.009] & 0.001 [0.009] & 0.002 [0.009] \\
\addlinespace
\multicolumn{5}{@{}l}{\textit{Panel D. $\ln M^*(10^{24})$ (truth 2.898)}} \\
Pooled Huber, $E$ free (ML practice) & 0.53 [1.61] & 0.55 [1.60] & $-$0.04 [1.41] & 0.64 [1.50] \\
Pooled NLS + generation FE & 0.01 [1.17] & 0.23 [0.73] & 0.10 [2.17] & 0.27 [0.80] \\
Lab FE & 0.00 [0.87] & 0.11 [0.65] & $-$0.06 [1.35] & 0.06 [0.83] \\
Lab $\times$ generation FE & 0.00 [0.37] & $-$0.01 [0.38] & $-$0.01 [0.33] & $-$0.02 [0.39] \\
ACF, current $c$ as instrument & 0.03 [1.47] & 0.10 [0.79] & 0.52 [2.83] & $-$0.03 [1.50] \\
ACF, lagged instruments only & 0.08 [1.58] & 0.02 [1.51] & 0.03 [1.50] & 0.03 [1.60] \\
IV (compute price) + path$^{\ddagger}$ & 1.27 [1.27] & 1.28 [1.28] & 1.32 [1.32] & 1.28 [1.28] \\
FOC system, $T$ ignored & 1.25 [1.25] & 1.25 [1.25] & 1.28 [1.28] & 1.25 [1.25] \\
FOC system, $T$ proxied & $-$0.02 [0.05] & $-$0.02 [0.05] & $-$0.03 [0.04] & $-$0.03 [0.05] \\
\bottomrule
\end{tabular}
\par\smallskip
\begin{tablenotes}Bias [root-mean-squared error] over 400 replications per compute rule. Simulated industry and compute rules: (a) exogenous budgets; (b) funding, $\ln C$ rises by $2x_{ft}$; (c) a loss target; (d) predetermined budgets, $\ln C$ rises by $2x_{f,t-1}$; setup in the text. $E$ is known to every estimator except the first. Both ACF rows use current within-lab-generation deviations of $n$ and $d$ as instruments, valid here by construction; without them the lagged moment set is under-identified. $^{\ddagger}$Median bias [median absolute error], because just-identified 2SLS has no finite moments; median first-stage $F$ is 17 to 18 under rules (a), (b), and (d) and 13 under (c). FOC system: the within-family loss equation stacked with the training-only first-order condition ($w=1$), ignoring or proxying inference demand $T$ (proxy with log noise of standard deviation 0.3). TFP growth: mean change in generation means of $y-\hat F(n,d)$. Every estimator is started from data-based values, none at the truth: nine random values for the pooled Huber fit, and nine combinations of $\alpha,\beta\in\{0.2,0.5,1\}$ with levels matched to mean output for the nonlinear least-squares estimators, keeping the lowest criterion; ACF starts from pooled NLS and the FOC systems from the lab $\times$ generation fit.\end{tablenotes}
\begin{tablenotes}[Source]Authors' simulations; replication code \texttt{code/analysis/m6\_montecarlo}.\end{tablenotes}
\end{table}


\begin{figure}[tp]
\centering
\includegraphics[width=\textwidth]{m6_montecarlo_industry_bias}
\caption{Bias of Industry Estimators by Compute Rule}\label{fig:app-mc-industry}
\begin{figurenotes}
Design~B, 400 replications per compute rule. Mean bias of $\hat\gamma$ (left), TFP growth per generation (center),
and $\ln\hat M^*(10^{24})$ (right) by estimator (rows) and compute rule (markers); the IV row reports median bias.
Mundlak and Heckman rows are close to lab fixed effects and pooled NLS, respectively, when there is no selection.
Truth: $\gamma=0.178$, TFP growth 0.15, $\ln M^*(10^{24})=2.898$.
\end{figurenotes}
\end{figure}

\subsection{Verification of the Transmission and Selection Formulas}\label{app:mc-formulas}

Table~\ref{tab:app-mc-transmission} checks the transmission-bias formula of Proposition~\ref{prop:transmission}
against simulation with $10^6$ draws per case. The simulated slopes match the closed form to within
$4\times10^{-4}$ in all twelve cases, on the expansion path and with lifetime-optimal allocations. The sign of the
bias follows the compute rule: exogenous budgets give $\gamma$, funding and predetermined budgets overstate it,
and a common loss target drives the slope to zero. The forward and reverse regressions bracket $\gamma$ under
exogenous budgets and dispersed targets, but under funding and predetermined budgets both exceed it. In samples of
240 lab-generations the nominal 95 percent OLS interval covers $\gamma$ in 95 percent of samples under exogenous
budgets, but in 53 percent under mild funding, at most 4 percent under stronger funding, and 16 percent under
predetermined budgets with $\lambda=1$; these coverage rates depend on the assumed exogenous dispersion of
compute. The residual dispersion also behaves as predicted: under funding with $\lambda=2$ the standard deviation
of the OLS residual is 0.229 against 0.255 for productivity plus noise, so cross-lab regressions understate TFP
dispersion.

Outcome-based release biases the compute slope toward zero. With exogenous budgets and an absolute release
threshold, the slope falls from 0.178 to 0.112, 0.082, and 0.057 as 25, 50, and 75 percent of runs are dropped;
combined with a funding rule ($\lambda=2$), the net bias moves from $+0.100$ without selection to $-0.006$ and
$-0.046$ with 50 and 75 percent dropped, so the net sign is ambiguous, as the theory implies. The allocation identities with factor-augmenting
productivity and the wedge identity of Proposition~\ref{prop:wedge} hold in the simulation to within $10^{-14}$.

\begin{table}[tp]
\centering
\caption{Transmission Bias in the Returns to Compute: Closed Form and Simulation}
\label{tab:app-mc-transmission}
\footnotesize
\setlength{\tabcolsep}{4pt}
\begin{tabular}{@{}lccccccc@{}}
\toprule
 & & \multicolumn{3}{c}{On the path ($T=0$), $n=10^6$} & $T>0$ & \multicolumn{2}{c}{$n=240$} \\
\cmidrule(lr){3-5}\cmidrule(lr){6-6}\cmidrule(l){7-8}
Compute rule & & Formula & Simulated & Reverse & Simulated & Mean & Coverage (\%) \\
\midrule
Exogenous &  & 0.1783 & 0.1781 & 0.542 & 0.1784 & 0.179 & 95 \\
Funding, $\lambda=$ & 0.5 & 0.2091 & 0.2095 & 0.512 & 0.2099 & 0.209 & 53 \\
Funding, $\lambda=$ & 1 & 0.2371 & 0.2370 & 0.481 & 0.2370 & 0.237 & 4 \\
Funding, $\lambda=$ & 2 & 0.2783 & 0.2782 & 0.429 & 0.2784 & 0.278 & 0 \\
Funding, $\lambda=$ & 4 & 0.3033 & 0.3034 & 0.359 & 0.3032 & 0.303 & 0 \\
Predetermined, $\lambda=$ & 1 & 0.2253 & 0.2254 & 0.487 & 0.2257 & 0.225 & 16 \\
Predetermined, $\lambda=$ & 2 & 0.2583 & 0.2582 & 0.435 & 0.2586 & 0.258 & 0 \\
Predetermined, $\lambda=$ & 4 & 0.2783 & 0.2783 & 0.359 & 0.2789 & 0.278 & 0 \\
Target, sd$(\bar y)=$ & 0 & 0.0000 & 0.0000 & -- & 0.0000 & 0.000 & 0 \\
Target, sd$(\bar y)=$ & 0.05 & 0.0069 & 0.0068 & 0.358 & 0.0068 & 0.007 & 0 \\
Target, sd$(\bar y)=$ & 0.1 & 0.0246 & 0.0246 & 0.223 & 0.0246 & 0.025 & 0 \\
Target, sd$(\bar y)=$ & 0.25 & 0.0891 & 0.0892 & 0.186 & 0.0886 & 0.089 & 0 \\
\bottomrule
\end{tabular}
\par\smallskip
\begin{tablenotes}Slope of $y=-\ln(L-E)$ on $c=\ln C$ across labs. Formula: $\gamma+\pi_1\mathrm{Var}(\omega)/(\pi_1^2\mathrm{Var}(\omega)+\mathrm{Var}(\eta))$ (Proposition~\ref{prop:A-transmission}) with $\gamma=0.1783$, sd$(\omega)=0.25$, sd$(\eta)=1$, persistence 0.8. Funding: $c=\lambda\omega+\eta$; predetermined: $c=\lambda\omega_{-1}+\eta$; target: labs buy the compute that reaches a target $\bar y$ with the stated dispersion (a common target gives slope 0). Simulated: OLS on $10^6$ draws, seed noise 0.05 (simulation standard errors at most 0.0003). Reverse: $\mathrm{Var}(y)/\mathrm{Cov}(c,y)$; ``--'' where $\mathrm{Cov}(c,y)=0$ and it is undefined. $T>0$: lifetime-optimal allocations with log-normal inference demand; under target rules an extra term of $-0.002$ to $-0.001$ appears because high-$T$ labs buy more compute. $n=240$: mean OLS slope and coverage of the nominal 95 percent OLS interval over 2,000 samples of 240 lab-generations; coverage rates are specific to sd$(\eta)=1$ and rise with more exogenous dispersion in compute.\end{tablenotes}
\begin{tablenotes}[Source]Authors' simulations; replication code \texttt{code/analysis/m6\_montecarlo}.\end{tablenotes}
\end{table}


\emph{Reproducibility.}---The code is in \texttt{code/analysis/m6\_montecarlo}. Every cell seeds its own random
number generator, so results do not depend on the number of processes; an independent rerun from raw inputs
reproduced every non-profile Design~A column and all original Design~B rows bit for bit. The full run takes about
45 minutes on 6 CPU processes.
